% Warnings
% chktex-file x
\section{Basic Concepts}

A secure system is :
\begin{itemize}
    \item \textbf{Confidential :} no attacker can \textbf{steal/read} data
    \item \textbf{Integrity : } no attacker can \textbf{modify} data/system behavior
    \item \textbf{Availability :} no attacker can disturb access 
\end{itemize}
We say that if a system is \textbf{vulnerable} it is exposed to \textbf{threats} that could lead to an \textbf{exploit} of the vulnerability so there is a \textbf{risk}.

\subsection{The STRIDE Threat Model}
Classification of security threats in six categories :
\begin{itemize}
    \item Spoofing : faking identity 
    \item Tampering : modification of smth 
    \item Repudiation : lack of logs to prove previous actions
    \item Information Disclosure : secrets get disclosed
    \item Denial of Service : make a service unavailable 
    \item Elevation of Privilege : getting higher privilege than initially
\end{itemize}

A voir si le reste du chapitre est vraiment intéressant pour l'examen ... 

We can do \textbf{access control} to restrict the access to certain resources to specific (groups of) users. We can implement it using both \textbf{authentication} (ie verifying the identity of the person) or \textbf{Authorization} (ie define who has access to what).

For files in Linux it is implemented using permissions defined on groups and users. For network it is implemented using firewalls.

As it is difficult to handle many privilege, it is recommended to use the \textbf{Principle of least privilege} : define the access control rules such that users only have access right needed to do their jobs, not more. It is the same for networking : no component should have more access than what is needed for its function (ie block everything and only allow what is needed).



